\pagebreak

\chapter{Setting up an FPGA board}\label{chap:boardsetup}
    
    The Gadget Factory papilio board provides a good example of
    the way 3PL is set up for a specific board.
    
    The following is a very simple 3PL application program for the papilio -
    
    \begin{lstlisting}[gobble=8, language=threepl]
        module "boards/gadget_factory/papilio_one_500.3pl"

        static(t, "log");

        seq {
            wait(0.5);
            t <- !t;
        }

        output(,, t, loc=A_pins[0]);
    \end{lstlisting}
    
    This trivial program outputs a 1Hz square wave on connector A pin 0 (P18).
    The only requirement in setting up the platform is the first statement
    which imports a file module specific to the papilio one 500.
    A 32MHz clock is established by default and is the current clock domain.
    
    The imported file looks as follows -

    \begin{lstlisting}[gobble=8, language=threepl]
        nosource

        // 3pl board support for GadgetFactory Papilio One 500k board

        part("XC3S500EVQ100-4");

        module "xilinx/xc3s.3pl"

        // all IO is 3.3V on this board
        str(lvcmos33, "LVCMOS33");

        // clock generator 32MHz
        iopins(c32_pin, "str", "P89", frequency=32.0e6, ibufg=true, iostandard=lvcmos33);

        // 48 I/O pins across A (wing 1A), B (wing 1B), C (wing 2A)
        iopins(global.A_pins, "[16]str", {
	        "P18", "P23", "P26", "P33", "P35", "P40", "P53", "P57",
	        "P60", "P62", "P65", "P67", "P70", "P79", "P84", "P86"
            }, iostandard=lvcmos33
        );
        iopins(global.B_pins, "[16]str", {
        .
        .
        .
        .

        iopins(global.flash_cs_pin, "str", "P24", iostandard=lvcmos33);
        iopins(global.flash_clk_pin, "str", "P50", iostandard=lvcmos33);
        iopins(global.flash_mosi_pin, "str", "P27", iostandard=lvcmos33);
        iopins(global.flash_miso_pin, "str", "P44", iostandard=lvcmos33);

        // FTDI USB channel B serial receive and transmit
        // Labelling is reversed to schematic!
        // 'tx' should be an FPGA output, 'rx' an input
        iopins(global.serial_rxd_pin, "str", "P88", iostandard=lvcmos33);
        iopins(global.serial_txd_pin, "str", "P90", iostandard=lvcmos33);

        // create 32MHz clock
        clock(c32_in, loc=c32_pin);
        simpleclock(global.c32, c32_in);
        useclock(c32);

        directive("comms_baud_rate", 115200); // define default serial device baud rate
        module "xilinx/uart_cpu.3pl"


        trailermodule ("create uart bus interface") {
            create_bus_interfaces();
        }
    \end{lstlisting}
    Some of the pin definitions have been omitted ... to shorten the listing.
    
    This file module specifies parameters for this particular board.
    
    The call to inbuilt procedure part() specifies the Xilinx FPGA part number.
    
    The next statement imports the file module xilinx/xc3s.3pl, described later.
    That file module specifies parameters appropriate for the particular FPGA.
    
    The calls to iopins() build a map of all input and output FPGA pins and
    attributes associated with them, such as the FPGA pin name and the I/O
    standard. Additional attributes can be specified when the pins are
    passed to an input(), output() or clock() inbuilt procedure.
    
    The first call to iopins() specifies the clock input pin (P89), the clock
    frequency, (32MHz), specifies that an IBUFG should be used and that the I/O
    standard is LVCMOS.
    
    A call to clock() creates a clock mode variable, \verb'c32_in'. The call to
    simpleclock() in xilinx/clocks.3pl creates a global clock mode variable c32,
    a clock manager (DCM) and a clock buffer. c32 is then made the current clock.
    
    Module \verb'xilinx/uart_cpu()' is imported to provide a CPU interface
    to the FPGA using the serial link. That interface is set up at the end
    by a trailer module which calls \verb'create_bus_interfaces()'
    
    Finally a trailer module is created to call \verb'create_bus_interfaces()'
    after the source code exits to build the logic required to flesh out
    the CPU/FPGA interface. Calls to the interface routines are made throughout
    the source code and this logic is required at the end to tie it all
    together.
    
    
    
    
    
    
    
    
    The file module xilinx/xc3s.3pl, while specific to the Spartan 3S, also
    imports libraries relevant to all Xilinx FPGAs.
    
    \begin{lstlisting}[gobble=8, language=threepl]
        // Include file for Xilinx Spartan 3 (XC3S) elements
        nosource

        presetdirective("postProcessTool", "ise");  // ISE will be used

        source "xilinx/common.3pl"      // xilinx common library
        module "xilinx/postprocess.3pl" // code to run ISE or Vivado
        module "stdlib.3pl"             // standard 3PL generic library
        module "xilinx/clocks.3pl"      // xilinx clock managers
        module "xilinx/io.3pl"          // xilinx I/O facilities

        // Specify allowed special I/O and clock elements valid for
        // this device.
        // Map "elementMap" is in common.3pl.
        elementMap["BUFG"] = true;
        elementMap["BUFGMUX"] = true;
        elementMap["BUFGCTRL"] = true;
        elementMap["DCM"] = true;

        leadermodule ("attributes", 100) {
            attributeprocedure(processPinAttributes);
        }

        trailermodule ("constraints", 100) {
            netlistcommentfromdirs();
            generateConstraints();
        }

        postmodule ("Xilinx postprocessing") {
            postprocess("xilinx/ise_fpga.cmd", directive("postProcessTool"), 9.1);
        }
    \end{lstlisting}

    This sets a directive to select the place-and-route tools to be used ("ise"
    or "vivado") and imports several appropriate libraries that might be needed.
    
    A leader module, which is executed prior to other source code, calls
    inbuilt procedure attributeprocedure() to notify 3PL that the library
    procedure processPinAttributes() in xilinx/common.3pl is to be called
    whenever an input() output() or clock() inbuilt procedure is called
    with I/O pins. This procedure is called for each pin and processes associated
    attributes and constraints for the pin.
    
    A trailer module, which is executed after the source program exits, calls two
    procedures, netlistcommentfromdirs() in stdlib.3pl and generateConstraints()
    in xilinx/common.3pl. The first of these simply copies all directives and
    their values as comments to the netlist (EDIF) file . The second generates
    constraints from information gathered during immediate execution.

    Finally a post-processing module is created which calls postprocess() in
    xilinx/postprocess.3pl.  (Note that post processing by the Xilinx tools can
    be supressed by the -s option when running 3pl. This might be done if the
    tools are to be run manually or perhaps on another computer. For example if
    3pl is executed on Apple OSX the Xilinx tools have to be executed on a Linux
    or Windows emulator such as Virtualbox or else run on another computer since
    Xilinx do not support OSX).




